
\begin{frame}{Notion de dépendance fonctionnelle}

\begin{defblock}{Dépendance fonctionnelle (DF)}
Il y a \alert{DF} $A \to B$ 
quand la connaissance de la valeur de $A$ implique la connaissance de la valeur de $B$.
\end{defblock}
\vfill 
Exemple\,: voici un schéma pour nos logements, \texttt{(code, nom, adresse, activité, description)},
avec les DF suivantes\,:

\vfill
\begin{redbullet}
  \item $code \to nom, adresse$
  \item $adresse \to code, nom$
 \item $code, activit\acute{e}  \to  description$
\end{redbullet}

\vfill

Les DF nous permettent de \alert{raisonner} sur le contenu d'une base et d'analyser sa qualité
\end{frame}


%%%%%%%%%%%%%%%
\begin{frame}{Les DF caractérisent les anomalies}
 
\begin{footnotesize}
\begin{tabular}{c|l|p{3.5cm}|c|p{4cm}}
\textbf{code} & \textbf{nom} & \textbf{adresse} & \textbf{activité}& \textbf{description} \\ \hline
   pi & U Pinzutu &  12 quai Napoléon & Plongée & Dans le golfe\\
   ta & Tabriz & Avenue du grand large & Plongée & Dans la rade\\
   pi & U Spuntinu &  12 avenue Paoli  & Voile & En club associatif\\
   ma & Macchia &  12 quai Napoléon & Gastronomie & Spécialités locales\\
\hline
\end{tabular}
\end{footnotesize}

\vfill
\begin{redbullet}
\item $code \to nom, adresse$ n'est pas respectée par les nuplets 1 et 3
    \item $adresse \to code, nom$ n'est pas respectée 
    par les nuplets 1 et 4
\item $code, activit\acute{e}  \to  description$ est toujours respectée
\end{redbullet}

\vfill

\alert{Remarque}\,: deux nuplets ayant les mêmes code logement 
et activité doivent être \alert{identiques} pour respecter les DF.
\end{frame}

%%%%%%%%%%%%%%%
\begin{frame}{Notion de clé}


\begin{defblock}{Clé}
Une clé d'une relation $R$ est un sous-ensemble \alert{minimal} $C$
des attributs tel que tout attribut de $R$ dépend fonctionnellement de $C$. 
\end{defblock}
\vfill

\begin{redbullet}
 \item $(code, activit\acute{e})$  est une clé\,: par transitivité, tous
 les autres attributs en dépendent
 \item $code$  n'est pas une clé, $activit\acute{e}$ et $description$ n'en dépendent pas
 \item $(code, activit\acute{e}, adresse)$  n'est pas une clé car non minimal
 
\end{redbullet}

\vfill

Deux nuplets avec la même valeur de clé doivent être identiques.
\end{frame}


%%%%%%%%%%%%%%%
\begin{frame}{Normalisation}
  
\begin{defblock}{Forme normale (Troisième, ou 3FN)}
  Un schéma de relation $R$ est en troisième forme normale quand
    tout attribut non-clé dépend directement d'une clé.
\end{defblock}
 
\vfill

\alert{Retenir} (petite approximation)\,: dans
   toute dépendance fonctionnelle $S \to A$ sur
   les attributs de $R$, $S$ est une clé. 

\vfill

Le schéma précédent n'est pas normalisé à cause des  DF
$code \to nom, adresse$ et $adresse \to code, nom$

\vfill

\alert{Intuition}\, une relation qui n'est pas en 3FN parle de deux
types d'objet différents, d'où redondances et incohérences.

\end{frame}


%%%%%%%%%%%%%%%
\begin{frame}{La bonne solution}
 
Voici le schéma \alert{normalisé} 
\vfill

\begin{redbullet}
  \item Logement (code, nom, adresse)
 \item Activité (code, activité, description)
\end{redbullet}

\vfill
\texttt{Logement} a deux clés, \texttt{code} et \texttt{adresse}

\vfill
\texttt{Activité} a une  clé, la paire  (\texttt{code, activité})

\vfill


\begin{redbullet}
  \item Ces relations sont  en troisième forme normale (3FN)
 \item On n'a pas perdu d'information car on peut reconstituer la relation initiale par jointure
\end{redbullet}

\end{frame}


%%%%%%%%%%%%%%%
\begin{frame}{Illustration}
 
\begin{tabular}{c|c|c}
\textbf{code} & \textbf{nom} & \textbf{adresse}\\ \hline
   pi & U Pinzutu &12 quai Napoléon\\
   ta & Tabriz  &Avenue du grand large\\
   ma & Macchia  &12 avenue Paoli\\
\hline
\end{tabular}

\vfill

\begin{tabular}{c|l|p{4cm}}
\textbf{code} & \textbf{activité} & \textbf{titre}\\ \hline
   pi& Plongée &  Dans le golfe\\
   pi & Voile&  En club associatif\\
   ta & Plongée&  Dans la rade\\
   ma & Gastronomie&  Spécialités locales\\
\hline
\end{tabular}

\vfill
\begin{defblock}{Contrainte d'unicité des clés}
L'unicité des valeurs de clé garantit le respect des DF.
\end{defblock}
\end{frame}

%%%%%%ù

\begin{frame}{À retenir}

La normalisation relationnelle décompose l'ensemble des attributs
en relations saines

\vfill
\begin{redbullet}
  \item Toute relation a une clé,
 \item Tous les attributs dépendent \alert{directement} de la clé (forme normale)
 \item Une valeur de clé doit apparaître une seule fois
 \item La décomposition en plusieurs relations est compensée par la possibilité
      d'effectuer des jointures
 \end{redbullet}

\vfill

\alert{Comment normalise-t-on\,?} À suivre.
\end{frame}

